\section*{How to use this companion}
\addcontentsline{toc}{section}{How to use this companion}
This companion and the article form a single proof package. The article gives
the objects, main arguments and conclusions; the entries below identify the
calculation behind each excursion and the point at which the main argument
resumes. Unqualified numbers refer to the article; lettered numbers refer
to this companion.

\begin{center}\small
\renewcommand{\arraystretch}{1.3}
\begin{tabularx}{\textwidth}{@{}>{\raggedright\arraybackslash}p{.24\textwidth}>{\raggedright\arraybackslash}X>{\raggedright\arraybackslash}p{.25\textwidth}@{}}
\toprule
Entry from the article & Calculation here & Return to the article\\\midrule
\Cref{lem:runge} & Coefficients, integrality and the growing-degree remainder in \cref{supp:lem:runge-calculation}. & Weighted defects, \cref{prop:defects}.\\
\Cref{macro:lem:pell,lem:split-bounded} & Uniform ideal and unit counts, then localization to two square classes in A.2--A.3. & Bounded components and the deep-core estimates.\\
\Cref{prop:cutoffs,thm:scalar-rates} & Weighted prime sums, degree bounds and scalar retention in B.1--B.3. & Scalar laws or the fields of Section~5.\\
\Cref{thm:word-field} & Marginal-cap substitutions and exponent bookkeeping in C.1. & Dictionary examples and consequences.\\
\Cref{thm:moving-comparison}(ii) & The two Stein bounds, Poisson filling and the single tail event in C.3--C.4. & The complete staircase and resolved kernels.\\
\Cref{thm:resolved-kernel} & Poisson atom estimates and the resolution cost in D.1--D.2. & Conditional laws in Section~6.\\
\Cref{thm:boundary} & Prime incidences, quotient rank and the post-quadratic bound in E.1--E.5. & The microscopic mass and the prefix law.\\
\Cref{thm:relative-bulk,thm:two-clock} & The normalized error table, finite-cylinder clock and affine conditioning in E.6--E.7. & The two-source comparison and its affine variant.\\\bottomrule
\end{tabularx}
\end{center}

\paragraph{Support conventions.}
The same symbols have different roles in different families of events.
For a run of length at least $L$, the system has $L$ relative equations
on $B=L+1$ vertices. An absolute word of length $B$ fixes $B$ values.
An exact run with excess $e$ uses $L+e+1$ relative equations on
$L+e+2$ vertices; recording its sign prescribes all these values.
After truncating at $e\le E$, the common maximal support has $Q+1$
vertices, where $Q=L+E+1$. In contrast, $Q_B=2^B$ is the exponential
parameter in the arithmetic relation profile.

\paragraph{Two different summations.}
For a spatially labelled field, marks are summed without weights before
the pair estimate, leaving the relation profile at the base length.
For an aggregated vector, the directional Stein weights are inserted
first, and the maximal-length profile is required. Section~C.3 keeps
this distinction explicit before any asymptotic absorption.
